% Topic 1.2.11 — Least Squares and the Gauss–Markov Theorem.

\section{Least Squares and the Gauss–Markov Theorem}
\label{sec:1.2.11-least-squares-gauss-markov}

\begin{ticket}
Least squares method. Point estimation of a vector parameter. System of normal equations. Properties of the least-squares estimator; Gauss–Markov theorem.
\end{ticket}

\subsection*{Theory}

We are given $n$ noisy observations of a linear function of an
unknown $p$-dimensional parameter and seek to estimate the parameter
by minimising the sum of squared residuals. The framework formalises
classical curve-fitting and is the prototype of every linear model
in statistics.

\begin{definition}[Linear regression model]
\label{def:linear-regression-model}
The \emph{linear regression model} (or \emph{linear measurement
model}) consists of:
\begin{itemize}
  \item a known \emph{design matrix} $\mat{X} \in \R^{n \times p}$
        of \emph{regressors} or covariates, with $n \geq p$;
  \item an unknown parameter vector $\vect{\beta} \in \R^p$;
  \item a vector $\vect{\varepsilon} \in \R^n$ of zero-mean,
        uncorrelated, equal-variance noise:
        $\E \vect{\varepsilon} = \vect{0}$ and
        $\Cov(\vect{\varepsilon}) = \sigma^2 \mat{I}_n$ for some
        unknown $\sigma^2 > 0$;
  \item observations $\vect{y} = \mat{X} \vect{\beta} + \vect{\varepsilon} \in \R^n$.
\end{itemize}
\end{definition}

\begin{definition}[Least-squares estimator]
\label{def:ols-estimator}
The \emph{ordinary least squares (OLS) estimator} of
$\vect{\beta}$ is any minimiser
\[
  \hat{\vect{\beta}} \;\in\; \argmin_{\vect{b} \in \R^p}\, \|\vect{y} - \mat{X} \vect{b}\|^2,
\]
where $\|\cdot\|$ is the Euclidean norm on $\R^n$.
\end{definition}

\begin{theorem}[Normal equations]
\label{thm:normal-equations}
Every OLS estimator $\hat{\vect{\beta}}$ satisfies the
\emph{normal equations}
\[
  \mat{X}\T \mat{X}\, \hat{\vect{\beta}} \;=\; \mat{X}\T \vect{y}.
\]
If $\mat{X}$ has full column rank $p$, then $\mat{X}\T \mat{X}$ is
invertible and the OLS estimator is unique:
\[
  \hat{\vect{\beta}} \;=\; (\mat{X}\T \mat{X})^{-1} \mat{X}\T \vect{y}.
\]
\end{theorem}
\proofof{normal-equations}

The geometric content of the normal equations: $\mat{X} \hat{\vect{\beta}}$
is the orthogonal projection of $\vect{y}$ onto the column space
$\im \mat{X} = \{\mat{X} \vect{b} : \vect{b} \in \R^p\}$, and the
residual $\vect{y} - \mat{X} \hat{\vect{\beta}}$ is orthogonal to
every column of~$\mat{X}$.

\begin{proposition}[Mean and covariance of the OLS estimator]
\label{prop:ols-properties}
Assume $\mat{X}$ has full column rank. Under
\cref{def:linear-regression-model}:
\begin{enumerate}[label=(\alph*)]
  \item $\hat{\vect{\beta}}$ is unbiased: $\E \hat{\vect{\beta}} = \vect{\beta}$;
  \item $\Cov(\hat{\vect{\beta}}) = \sigma^2 (\mat{X}\T \mat{X})^{-1}$;
  \item the residual sum of squares
        $\mathrm{RSS} = \|\vect{y} - \mat{X} \hat{\vect{\beta}}\|^2$
        satisfies $\E \mathrm{RSS} = (n - p) \sigma^2$.
\end{enumerate}
\end{proposition}
\proofof{ols-properties}

The estimator $\hat{\sigma}^2 = \mathrm{RSS}/(n - p)$ is therefore an
unbiased estimator of the noise variance.

\begin{theorem}[Gauss–Markov]
\label{thm:gauss-markov}
Assume \cref{def:linear-regression-model} with $\mat{X}$ of full
column rank. Among all \emph{linear unbiased estimators} of
$\vect{\beta}$ — that is, estimators of the form $\tilde{\vect{\beta}} = \mat{C} \vect{y}$
with $\mat{C} \in \R^{p \times n}$ and $\E \tilde{\vect{\beta}} = \vect{\beta}$
identically in $\vect{\beta}$ — the OLS estimator $\hat{\vect{\beta}}$
has minimum variance, in the sense that
\[
  \Cov(\tilde{\vect{\beta}}) - \Cov(\hat{\vect{\beta}}) \quad \text{is positive semi-definite.}
\]
Equivalently, for every linear functional $\vect{a}\T \vect{\beta}$
of the parameter, $\vect{a}\T \hat{\vect{\beta}}$ is the
\emph{best linear unbiased estimator} (BLUE) of $\vect{a}\T \vect{\beta}$:
$\Var(\vect{a}\T \tilde{\vect{\beta}}) \geq \Var(\vect{a}\T \hat{\vect{\beta}})$
for every linear unbiased $\tilde{\vect{\beta}}$.
\end{theorem}
\proofof{gauss-markov}

\begin{remark}
Gauss–Markov assumes only zero mean, equal variance, and lack of
correlation of the noise — \emph{not} normality. If the noise is
additionally Gaussian, $\hat{\vect{\beta}}$ is also the maximum
likelihood estimator and attains the Cramér–Rao lower bound
\cite[Ch.~VII, \S~6]{Shiryaev1989}; in that case it is BLUE among
\emph{all} unbiased estimators (linear or not). If the noise has
non-trivial covariance $\Cov(\vect{\varepsilon}) = \mat{\Sigma}$
with known $\mat{\Sigma} \succ 0$, the BLUE is the
\emph{generalised} least-squares estimator
$\hat{\vect{\beta}}_{\mathrm{GLS}} = (\mat{X}\T \mat{\Sigma}^{-1} \mat{X})^{-1} \mat{X}\T \mat{\Sigma}^{-1} \vect{y}$.
\end{remark}

\subsection*{Proofs}

\begin{proof}[\Cref{thm:normal-equations}]
\proofanchor{normal-equations}
The squared-error objective $L(\vect{b}) = \|\vect{y} - \mat{X} \vect{b}\|^2$
expands as
$L(\vect{b}) = \vect{y}\T \vect{y} - 2 \vect{b}\T \mat{X}\T \vect{y} + \vect{b}\T \mat{X}\T \mat{X}\, \vect{b}$.
This is a convex quadratic in~$\vect{b}$, so the gradient
\[
  \nabla L(\vect{b}) = -2\, \mat{X}\T \vect{y} + 2\, \mat{X}\T \mat{X}\, \vect{b}
\]
vanishes precisely at the minimum (\cref{thm:cont-partials-diff}
applied to the convex case is a sufficient condition for the
critical point being global). Setting $\nabla L(\hat{\vect{\beta}}) = \vect{0}$
gives $\mat{X}\T \mat{X}\, \hat{\vect{\beta}} = \mat{X}\T \vect{y}$.
If $\mat{X}$ has full column rank, then $\mat{X}\T \mat{X}$ is
invertible (it is a $p \times p$ Gram matrix of linearly independent
columns; cf.\ \cref{sec:1.1.07-quadratic-forms}), so the solution
is the unique $(\mat{X}\T \mat{X})^{-1} \mat{X}\T \vect{y}$.
\end{proof}

\begin{proof}[\Cref{prop:ols-properties}]
\proofanchor{ols-properties}
Let $\mat{H} = \mat{X} (\mat{X}\T \mat{X})^{-1} \mat{X}\T$ — the
orthogonal projector onto $\im \mat{X}$. Note $\mat{H}$ is symmetric
and idempotent: $\mat{H}\T = \mat{H}$, $\mat{H}^2 = \mat{H}$.

\emph{(a)} Substitute $\vect{y} = \mat{X} \vect{\beta} + \vect{\varepsilon}$:
$\hat{\vect{\beta}} = (\mat{X}\T \mat{X})^{-1} \mat{X}\T (\mat{X} \vect{\beta} + \vect{\varepsilon}) = \vect{\beta} + (\mat{X}\T \mat{X})^{-1} \mat{X}\T \vect{\varepsilon}$.
Take expectation: $\E \hat{\vect{\beta}} = \vect{\beta} + (\mat{X}\T \mat{X})^{-1} \mat{X}\T \E \vect{\varepsilon} = \vect{\beta}$.

\emph{(b)} From the same decomposition,
$\Cov(\hat{\vect{\beta}}) = (\mat{X}\T \mat{X})^{-1} \mat{X}\T \Cov(\vect{\varepsilon}) \mat{X} (\mat{X}\T \mat{X})^{-1} = \sigma^2 (\mat{X}\T \mat{X})^{-1}$
using $\Cov(\vect{\varepsilon}) = \sigma^2 \mat{I}_n$.

\emph{(c)} The residual is
$\vect{y} - \mat{X} \hat{\vect{\beta}} = (\mat{I}_n - \mat{H}) \vect{y} = (\mat{I}_n - \mat{H}) \vect{\varepsilon}$
(since $(\mat{I} - \mat{H}) \mat{X} = \mat{0}$). Thus
$\mathrm{RSS} = \vect{\varepsilon}\T (\mat{I}_n - \mat{H}) \vect{\varepsilon}$,
a quadratic form in zero-mean noise. Taking expectation,
$\E \mathrm{RSS} = \tr\bigl((\mat{I}_n - \mat{H}) \Cov(\vect{\varepsilon})\bigr) = \sigma^2 \tr(\mat{I}_n - \mat{H}) = \sigma^2 (n - \rank \mat{H}) = \sigma^2 (n - p)$,
using the standard identity $\E[\vect{\varepsilon}\T \mat{A} \vect{\varepsilon}] = \tr(\mat{A} \Cov(\vect{\varepsilon}))$
for zero-mean~$\vect{\varepsilon}$ and the fact that
$\rank \mat{H} = \rank \mat{X} = p$.
\end{proof}

\begin{proof}[\Cref{thm:gauss-markov}]
\proofanchor{gauss-markov}
Let $\tilde{\vect{\beta}} = \mat{C} \vect{y}$ be linear and unbiased.
The constraint
$\E \tilde{\vect{\beta}} = \mat{C} \mat{X} \vect{\beta} = \vect{\beta}$
for every $\vect{\beta}$ forces $\mat{C} \mat{X} = \mat{I}_p$. Write
$\mat{C} = (\mat{X}\T \mat{X})^{-1} \mat{X}\T + \mat{D}$ where
$\mat{D} \in \R^{p \times n}$. Then $\mat{D} \mat{X} = \mat{0}_p$ from
the constraint.

The covariance is
\[
  \Cov(\tilde{\vect{\beta}}) = \mat{C}\, \Cov(\vect{\varepsilon})\, \mat{C}\T = \sigma^2\, \mat{C} \mat{C}\T.
\]
Expand $\mat{C} \mat{C}\T = \bigl((\mat{X}\T \mat{X})^{-1} \mat{X}\T + \mat{D}\bigr)\bigl(\mat{X} (\mat{X}\T \mat{X})^{-1} + \mat{D}\T\bigr)$.
The cross terms vanish by $\mat{D} \mat{X} = \mat{0}$:
\[
  \mat{C} \mat{C}\T = (\mat{X}\T \mat{X})^{-1} + \mat{D} \mat{D}\T.
\]
Hence $\Cov(\tilde{\vect{\beta}}) - \Cov(\hat{\vect{\beta}}) = \sigma^2 \mat{D} \mat{D}\T$,
which is positive semi-definite (a Gram matrix). Equality holds iff
$\mat{D} = \mat{0}$, i.e.\ $\tilde{\vect{\beta}} = \hat{\vect{\beta}}$.

The linear-functional version follows: for any
$\vect{a} \in \R^p$,
$\Var(\vect{a}\T \tilde{\vect{\beta}}) - \Var(\vect{a}\T \hat{\vect{\beta}}) = \sigma^2 \vect{a}\T \mat{D} \mat{D}\T \vect{a} = \sigma^2 \|\mat{D}\T \vect{a}\|^2 \geq 0$.
\end{proof}

\subsection*{Examples}

\begin{example}[Simple linear regression]
\label{ex:simple-linear-regression}
Fit $y_i = a + b x_i + \varepsilon_i$ for $i = 1, \dots, n$ — the
$p = 2$ case with
$\mat{X} = \begin{pmatrix} 1 & x_1 \\ \vdots & \vdots \\ 1 & x_n \end{pmatrix}$.
Setting $\bar x = \tfrac{1}{n}\sum x_i$ and $\bar y = \tfrac{1}{n}\sum y_i$,
the normal equations
$\mat{X}\T \mat{X} \begin{pmatrix} \hat a \\ \hat b \end{pmatrix} = \mat{X}\T \vect{y}$
solve to
\[
  \hat b = \frac{\sum_{i} (x_i - \bar x)(y_i - \bar y)}{\sum_{i} (x_i - \bar x)^2}, \qquad
  \hat a = \bar y - \hat b\, \bar x.
\]
By \cref{prop:ols-properties}\,(b),
$\Var(\hat b) = \sigma^2 / \sum_i (x_i - \bar x)^2$ — the precision
improves with the spread of the $x_i$.
\end{example}

\begin{example}[Numerical fit]
\label{ex:numerical-fit}
Data: $(x_i, y_i) = (1, 1), (2, 3), (3, 4), (4, 6), (5, 5)$.
Compute $\bar x = 3$, $\bar y = 19/5 = 3.8$,
$\sum (x_i - \bar x)(y_i - \bar y) = (-2)(-2.8) + (-1)(-0.8) + 0 + (1)(2.2) + (2)(1.2) = 5.6 + 0.8 + 0 + 2.2 + 2.4 = 11.0$,
$\sum (x_i - \bar x)^2 = 4 + 1 + 0 + 1 + 4 = 10$. Hence $\hat b = 1.1$
and $\hat a = 3.8 - 1.1 \cdot 3 = 0.5$. The fitted line is
$\hat y = 0.5 + 1.1 x$.
\end{example}

\begin{problem}
\label{prob:gauss-markov-comparison}
For the simple linear regression of \cref{ex:simple-linear-regression}
with $n = 3$ observations at $x_1 = -1, x_2 = 0, x_3 = 1$, consider
the two competing estimators of the slope:
\begin{itemize}
  \item the OLS estimator $\hat b$ from the normal equations;
  \item the naive ``two-point'' estimator $\tilde b = (y_3 - y_1)/(x_3 - x_1) = (y_3 - y_1)/2$.
\end{itemize}
Show both are linear and unbiased and compare their variances.
\end{problem}

\begin{proof}[Solution]
With $\bar x = 0$, the OLS slope simplifies to
$\hat b = \sum_i x_i (y_i - \bar y)/\sum_i x_i^2 = ((-1) y_1 + 0 \cdot y_2 + 1 \cdot y_3)/2 - \bar y \cdot 0 = (y_3 - y_1)/2$.
So in this small symmetric design, the two estimators \emph{coincide}.
$\hat b = \tilde b$ are equal as linear combinations of $\vect{y}$,
hence have the same mean and variance — Gauss–Markov gives nothing
to choose between them. The story changes if we add a fourth
observation at $x_4 = 0$ which the two-point estimator simply
ignores; OLS uses it (slightly improving precision) by reducing
$\Var(\bar y)$, hence reducing the residual that affects the slope
estimate. This is the practical content of Gauss–Markov: OLS uses
\emph{all} information optimally under the linear-unbiased
constraint.
\end{proof}
